\documentclass[../DoAn.tex]{subfiles}
\begin{document}
\label{ch:mo_dau}
% =============================================================
% CHƯƠNG 1 - GIỚI THIỆU ĐỀ TÀI
% =============================================================

\section{Đặt vấn đề}
\label{sec:dat-van-de}

Trong vòng hai thập kỷ qua, sự bùng nổ của điện toán đám mây, Internet vạn vật (IoT) và các dịch vụ thời gian thực đã tạo nên áp lực rất lớn lên hạ tầng mạng truyền thống. Mỗi dịch vụ có yêu cầu khác biệt về băng thông, độ trễ và đặc tính bảo mật, dẫn đến nhu cầu chạy nhiều mạng logic độc lập trên cùng một hạ tầng vật lý. Công nghệ \textbf{ảo hoá mạng} (Network Virtualization) ra đời nhằm đáp ứng nhu cầu này, đóng vai trò nền tảng cho các kiến trúc \textit{Software-Defined Networking} (SDN) và \textit{Network Function Virtualization} (NFV) hiện đại\cite{peterson2007computer,kreutz2015sdn_survey,mijumbi2016nfv_survey}.

Trong mô hình ảo hoá mạng, nhà cung cấp hạ tầng (Infrastructure Provider - InP) đồng thời phục vụ nhiều nhà cung cấp dịch vụ (Service Provider - SP). Mỗi SP gửi tới InP một \textit{yêu cầu mạng ảo} (Virtual Network Request - VNR) gồm tô-pô nút - cạnh cùng các yêu cầu tài nguyên (CPU cho nút, băng thông cho cạnh). Nhiệm vụ của InP là tìm cách \textbf{nhúng} từng VNR lên mạng vật lý sao cho thoả mãn ràng buộc tài nguyên và tối ưu một hàm mục tiêu đã định, chẳng hạn tỉ lệ chấp nhận yêu cầu hay tỉ số doanh thu trên chi phí (R/C)\cite{fischer2013vne_survey,cao2019vne_survey}.

Bài toán nhúng mạng ảo là cốt lõi của các kịch bản 5G/6G network slicing, multi-tenant cloud và edge computing. Hiệu quả của thuật toán VNE ảnh hưởng trực tiếp tới khả năng tận dụng hạ tầng, chất lượng dịch vụ và doanh thu của nhà cung cấp. Mặt khác, VNE đã được chứng minh là NP-Hard~\cite{chowdhury2009vine}: ngay cả phiên bản ánh xạ nút đơn lẻ đã quy được về bài toán Multi-Way Separator Problem, do vậy không thể giải tối ưu trong thời gian đa thức ở quy mô thực tế. Tính chất NP-Hard, kết hợp với tô-pô mạng vật lý có thể lên đến hàng trăm nút và VNR đến liên tục theo thời gian, đặt ra một thách thức tính toán rất lớn cho thiết kế thuật toán.

Trong các kịch bản thực tế, bài toán còn phức tạp hơn do tính \textbf{trực tuyến}: VNR đến theo phân phối Poisson, có thời gian sống ngẫu nhiên, và quyết định nhúng cho từng VNR phải đưa ra ngay lập tức mà không biết trước các yêu cầu trong tương lai. Một thuật toán tốt phải vừa khai thác hiệu quả tài nguyên hiện tại, vừa giữ ``khoảng trống'' chiến lược để phục vụ các yêu cầu sau.

Các hướng nghiên cứu hiện nay gồm ba nhóm chính: (i) heuristic/metaheuristic dựa trên tri thức chuyên gia (nhanh, đơn giản nhưng khó tổng quát và dễ rơi vào tối ưu cục bộ), (ii) lập trình toán học ILP (đảm bảo tối ưu nhưng chỉ áp dụng cho bài toán quy mô nhỏ), và (iii) học tăng cường sâu (học chính sách end-to-end nhưng phụ thuộc reward shaping, chưa khai thác tốt tri thức tô-pô và khó hội tụ trong không gian hành động lớn); phần khảo sát chi tiết được trình bày ở Chương~\ref{ch:bai_toan}. Những hạn chế của hướng học tăng cường là động lực chính cho phương pháp đề xuất trong đồ án.

\section{Mục tiêu của đồ án}

Trên cơ sở phân tích các hạn chế ở trên, đồ án đặt ra ba mục tiêu cụ thể:

\begin{enumerate}
\item \textbf{Trích xuất biểu diễn tô-pô đầy đủ} cho cả mạng vật lý và yêu cầu mạng ảo bằng mạng nơ-ron đồ thị (GNN), thay cho các đặc trưng thủ công của heuristic.
\item \textbf{Thu hẹp không gian hành động} bằng cơ chế attention chọn tập ứng viên (candidate) cho từng nút ảo, giúp tác nhân tập trung học sự đánh đổi giữa các phương án thực sự khả thi.
\item \textbf{Tận dụng tri thức từ heuristic} thông qua giai đoạn tiền huấn luyện có giám sát dựa trên lời giải của heuristic MP-VNE, rồi tinh chỉnh chính sách bằng học tăng cường để vượt qua trần hiệu năng của heuristic.
\end{enumerate}

Định hướng giải pháp tổng thể là xây dựng kiến trúc \textit{CARL-VNE} (Candidate-Anchored Reinforcement Learning for Virtual Network Embedding, học tăng cường chọn ứng viên có neo) gồm một bộ mã hoá đồ thị \textit{phân cấp đa miền} (GCN nội miền kết hợp GCN liên miền), một đầu chọn ứng viên cross-attention cùng một bộ phê bình giá trị, kết hợp với quy trình huấn luyện hai giai đoạn: tiền huấn luyện có giám sát rồi tinh chỉnh bằng phương pháp \textit{actor-critic} (gradient chính sách với phê bình giá trị làm baseline) có neo Kullback-Leibler. Một phát hiện then chốt là \textit{lựa chọn ứng viên} mới là đòn bẩy hiệu quả của học tăng cường, nên toàn bộ tín hiệu học được dồn vào đầu chọn ứng viên. Kết quả là một tác nhân vừa kế thừa các quy luật hợp lý của heuristic, vừa có khả năng tự khám phá những phương án ánh xạ tốt hơn mà heuristic không nhìn thấy.

\section{Đóng góp của đồ án}

Đồ án có ba đóng góp chính sau:
\begin{enumerate}
\item Đề xuất kiến trúc \textit{CARL-VNE} với bộ mã hoá đồ thị phân cấp đa miền (GCN nội miền + liên miền) kết hợp đầu chọn ứng viên cross-attention tích hợp đặc trưng chi phí tường minh (cost-aware) và mặt nạ khả thi cứng cho ràng buộc CPU.
\item Đề xuất quy trình huấn luyện hai giai đoạn gồm tiền huấn luyện có giám sát từ lời giải của heuristic MP-VNE và tinh chỉnh bằng phương pháp actor-critic (phê bình giá trị làm baseline) có neo Kullback-Leibler, cùng phát hiện rằng lựa chọn ứng viên là đòn bẩy hiệu quả và việc làm cho nó ngẫu nhiên trong huấn luyện là điều kiện cần để vượt trần của mô hình bắt chước.
\item Xây dựng bộ benchmark đa miền $50$ nút (chia thành $10$ miền) với $21$ bộ dữ liệu phân phối yêu cầu mạng ảo, và tiến hành đánh giá so sánh có hệ thống với baseline heuristic MP-VNE.
\end{enumerate}

\section{Bố cục của đồ án}

Phần còn lại của báo cáo được tổ chức như sau. Chương~\ref{ch:nen_tang} trình bày cơ sở lý thuyết, bao gồm các kiến thức về học tăng cường sâu, mạng nơ-ron đồ thị và cơ chế attention. Chương~\ref{ch:bai_toan} khảo sát các hướng tiếp cận hiện có và hình thức hoá bài toán nhúng mạng ảo đa miền: mô hình mạng vật lý đa miền và mạng ảo, hai pha ánh xạ nút và liên kết, hàm mục tiêu, các độ đo hiệu năng và phát biểu dưới dạng quy hoạch nguyên tuyến tính. Chương~\ref{ch:de_xuat} trình bày phương pháp đề xuất của đồ án, gồm kiến trúc tổng thể, bộ mã hoá đồ thị, bộ giải mã chọn ứng viên, môi trường học tăng cường, quy trình huấn luyện hai giai đoạn và phân tích lý thuyết về độ phức tạp cùng đặc tính hội tụ. Chương~\ref{ch:thuc_nghiem} mô tả thiết lập thực nghiệm, các tập dữ liệu, độ đo so sánh và phân tích kết quả. Phần Kết luận tổng kết các kết quả đạt được, các hạn chế còn tồn tại và đề xuất các hướng nghiên cứu tiếp theo.

\end{document}
